\section{Model}\label{sec:model}

% Formal skeleton assembled from pipeline/dynamic_game/PROOF_NOTES.md.
% Connective prose is deliberately minimal and factual; TODO(BH) markers
% flag where argument-carrying prose is to be written.

\subsection{Environment}\label{sec:env}

Time is discrete, $t = 0, 1, \ldots, T$ (quarters). There are $n$
protocols indexed by $i$, protocol $i$ holding value (TVL) $V_{it} > 0$,
and a coverage pool with liquidity-provider (LP) capital $S_t \geq 0$.
Each period, protocol $i$ chooses a security level $h_i \geq 0$ at flow
cost $c_h h_i V_{it}$ and posts collateral $C_i \in [0, \bar{c}\,V_{it}]$
into the pool. Total pool capital is $A_t = S_t + \sum_i C_i$.

\paragraph{Attacker.} A profit-motivated attacker observes
$(V_{it}, h_i)$ and chooses effort $e \geq 0$ per target. Extraction
technology and cost are
\begin{equation}\label{eq:attacker-tech}
  \blast(e, h) = \bigl(1 - e^{-a e}\bigr)(1 + h)^{-\zeta},
  \qquad
  c(e, V) = \bigl(F_0 + \gamma_c e^{\kappa}\bigr) V^{\eta},
\end{equation}
so that a successful exploit extracts $\varepsilon \blast(e,h) V$ with
execution noise $\varepsilon \in [0,1]$,
$\E[\varepsilon] = \bar\varepsilon$. The elasticity $\eta$ of attack cost
with respect to target value is the paper's key structural parameter
(\cref{sec:calibration}); $\kappa$ governs the effort cost curvature.
The attacker's rent against target $(V, h)$ and the induced hack
probability are
\begin{equation}\label{eq:rent}
  \rent(V, h) = \max_{e \geq 0}\;
    \bar\varepsilon\, \blast(e, h)\, V - c(e, V),
  \qquad
  p(V, h) = q \cdot \Bigl(1 - e^{-\rent(V,h)_+ / \bar{F}}\Bigr),
\end{equation}
with $q$ the per-quarter opportunity arrival rate and $\bar{F}$ the
scale of the (smoothed) participation margin. % TODO(BH): footnote on the
% smoothing (heterogeneous outside options) and the exact-threshold variant.

\paragraph{Mechanism.} The pool grants protocol $i$ coverage
\begin{equation}\label{eq:coverage}
  \cov_i = s\,\mu\, C_i^{\theta}\,\bigl(1 + \xi(h_i)\bigr),
  \qquad \theta \in (0,1),
\end{equation}
where $\xi(\cdot)$ is an increasing security kicker and
$s = \min\{1, U_{\max}/U\}$ a prudential cap on utilization
$U = \sum_i \cov_i / A$. Pool capital earns
$\poolret(A) \geq \rmkt$, with either $\poolret$ fixed (baseline) or
declining in pool size toward the market rate,
$\poolret(A) = \rmkt + (\poolret_0 - \rmkt)\,K/(K + A)$ (the
endogenous-$\alpha$ extension). Revenue is split between LPs and
protocols by a utilization- and risk-dependent share $\gamma$;
protocols' share is pro-rata in posted collateral. On a hack with loss
$L_i = \varepsilon \blast V_{it}$, the pool pays the claim
\begin{equation}\label{eq:claim}
  \text{claim}_i = (1 - d)\,\min\{\cov_i,\, L_i\},
\end{equation}
where $d \in [0,1]$ is a coinsurance rate ($d = 0$ in the baseline
mechanism), and --- the design's distinguishing instrument --- the hacked
protocol's collateral $C_i$ is \emph{forfeited} to the pool.
% TODO(BH): institutional paragraph — parametric on-chain trigger, no
% discretionary negligence adjudication; map to MARBLE mechanism Eqs. 2–8.

\paragraph{Protocol payoff.} Protocol $i$'s per-quarter objective is
\begin{equation}\label{eq:protocol-utility}
\begin{aligned}
  U_i(C_i, h_i) ={}& p_i\,\E\bigl[(1-d)\min\{\cov_i, L_i\}\bigr]
     \;-\; p_i\, C_i \,\ind{\text{forfeiture}}
     \;+\; (1 - p_i)\, y_i(C_i) \\
  &- \tfrac{\rmkt}{4} C_i
     \;-\; \tfrac{c_h}{4}\, h_i V_i
     \;-\; p_i \E[L_i]
     \;-\; \rho_i\, p_i\,\bigl(\E[L_i] - \E[\text{claim}_i]\bigr),
\end{aligned}
\end{equation}
where $y_i(C_i)$ is the pro-rata yield share and $\rho_i$ a
protocol-specific loss-aversion weight on uninsured losses.
% TODO(BH): align notation table with MARBLE Eq. (7) and the addon README.

\subsection{Timing and solution concept}\label{sec:timing}

Within each quarter: (1) the mechanism state
$(A, \gamma, s, \text{hazard})$ is quoted; (2) protocols simultaneously
choose $(C_i, h_i)$; (3) the attacker best-responds target-by-target;
(4) hacks realize with noise $\varepsilon$; (5) settlement: claims are
paid, collateral of hacked protocols is forfeited to the pool, LP capital
adjusts, and TVL transitions. % TODO(BH): timing figure (pipeline
% fig_game_structure) and transition equations.

\begin{definition}[Aggregate-taking stage equilibrium]\label{def:ate}
A profile $\{(C_i^*, h_i^*)\}_i$ together with aggregates
$(\gamma^*, s^*, y^*)$ is an aggregate-taking stage equilibrium if
(i) each $(C_i^*, h_i^*)$ maximizes \eqref{eq:protocol-utility} given the
aggregates and the attacker's best response \eqref{eq:rent}, and (ii) the
aggregates are generated by the profile.
\end{definition}

% TODO(BH): remark on epsilon-equilibria in the finite game and the
% continuum limit (existence exact); cite PROOF_NOTES Steps 3 and the
% certification table.

\subsection{Equilibrium existence}\label{sec:existence}

\begin{lemma}[Own-collateral concavity]\label{lem:concavity}
Suppose pool revenue $R(A) = \poolret(A)\,A$ is concave with
$R(0) = 0$. Then protocol $i$'s objective
\eqref{eq:protocol-utility} is concave in own collateral $C_i$ on
$[0, \bar{c} V_i]$ for every fixed $h_i$.
\end{lemma}

\begin{theorem}[Existence]\label{thm:existence}
Under \cref{lem:concavity} and continuity of the aggregate maps, an
aggregate-taking stage equilibrium in pure strategies exists in the
continuous-collateral game with finite security grid; in the
continuum-of-protocols limit an exact equilibrium exists with
$h$ continuous. % TODO(BH/formal): the continuum statement is the target
% theorem for this paper; the finite-grid game admits epsilon-equilibria
% with certified bounds (Appendix~\ref{app:proofs}).
\end{theorem}

\subsection{The carry trade as a security subsidy}\label{sec:carry}

\begin{proposition}[Advantageous selection and carry--security
complementarity]\label{prop:carry}
Let $\chi_i = (1 - p_i)\,\partial y_i/\partial C_i - \rmkt/4 -
p_i \ind{\text{forfeiture}}$ denote the carry margin on posted
collateral. Then (i) $\partial \chi_i / \partial p_i < 0$: safer
protocols earn a strictly higher carry margin, so collateral posting
selects advantageously on security; (ii) with forfeiture on, the cross
partial of $U_i$ in $(C_i, h_i)$ is positive wherever the deterrence
channel $-\partial p_i/\partial h_i\,(\partial y_i/\partial C_i + 1) > 0$
dominates: collateral and security are complements.
\end{proposition}

\subsection{Forfeiture decouples deterrence from indemnity}\label{sec:decoupling}

Let \emph{family A} (indemnity-only contracts) be the set of mechanisms
with forfeiture off and coinsurance $d \in [0,1]$; every classical
indemnity-side instrument (deductible, coinsurance, coverage denial,
negligence exclusion) is a state-contingent choice of $d$. Let
$s(d) \in [0,1]$ be the equilibrium indemnified share of realized losses
and $h^*(d)$ the equilibrium security level.

\begin{lemma}[Indemnity irrelevance to the attacker]\label{lem:irrelevance}
Fix the defender profile $\{(C_i, h_i)\}_i$. The attacker's problem
\eqref{eq:rent} contains no indemnity-side object: $e^*$, $\rent$, and
$p$ are independent of $d$ and of the claim rule. Hence insurance design
affects hack arrivals only through the induced defender choices
$(C_i, h_i)$.
\end{lemma}

\begin{lemma}[The indemnity frontier]\label{lem:frontier}
Within family A, the on-hack utility differential that disciplines
$h_i$ is bounded by the withheld claim,
$\Delta_A \leq \E[\min\{\cov_i, L_i\}]$, so discipline requires
withholding transfer: over the moral-hazard region the attainable set of
family A traces a decreasing frontier $s \mapsto h^*(s)$, and by
\cref{lem:irrelevance} hack intensity increases along it.
\end{lemma}

\begin{proposition}[Decoupling]\label{prop:decoupling}
The forfeiture mechanism ($d = 0$, forfeiture on) attains a
transfer--security pair outside the family-A attainable set:
(i) no indemnity-only contract achieves both the mechanism's indemnified
share and its security level simultaneously; (ii) the mechanism's
equilibrium security strictly exceeds the family-A frontier evaluated at
the mechanism's transfer level. The mechanism does not maximize the
indemnified share itself: concave coverage caps tail claims by design,
so the escape is in the pair, not in either coordinate alone.
(iii) The value of the second instrument is amplified by the strategic
attacker: shutting down the attacker's extensive-margin response
compresses the security gap between the forfeiture mechanism and the
best family-A contract.
% Certified numerically on the calibrated game (strict criteria, PASS):
% frontier Spearman rho = -1.00 (p < 1e-4); at the mechanism's transfer
% share s = 0.876 the interpolated family-A frontier gives h = 0.40
% (~8 hacks/yr); the mechanism delivers h* = 1.15 (the no-insurance
% security level) at 2.1 hacks/yr -- escape gap 0.75 in h. See
% pipeline/dynamic_game/outputs/results_decoupling.json.
% TODO(BH/formal): analytical version for the continuum model --
% supermodularity of the payout term in (-d, h) where coverage binds;
% part (iii) comparative static in d(beta)/dh and the participation
% elasticity via F-bar.
\end{proposition}

\begin{corollary}[Why the hostage is postable]\label{cor:hostage}
The classical objections to forfeitable bonds --- agents cannot or will
not post them, and the bond-holder may seize them opportunistically ---
are answered inside the mechanism: the carry margin
$\chi_i > 0$ (\cref{prop:carry}) compensates the expected forfeiture, and
the parametric on-chain trigger removes adjudicator discretion.
% TODO(BH): connect to Williamson (1983), Dickens et al. (1989) in the
% positioning section; this corollary carries the "why blockchain"
% argument in contract-theoretic form.
\end{corollary}
